The rest of this section consists of the proof of Theorem~\ref{thm-PPconj}.

\subsection{Reductions} We start by reducing Theorem~\ref{thm-PPconj} to simpler cases. First notice that all expressions in Theorem~\ref{thm-PPconj} are defined over the field $\FF_2$. Hence we may assume that $k=\FF_2$. 

Recall that we wrote $W_\Delta = \sum_i W_{\Pi_i}$ in Theorem~\ref{thm-W-sum}. 

\begin{lemma} \label{lem-redPi}
Theorem~\ref{thm-PPconj} for all $\Pi_i$ implies it for $\Delta$.
\end{lemma}

\begin{proof}
This follows directly from the statement of the theorem using the characteristic $2$ assumption and the  observation that if a monomial $x_Ix_J$ restricts to a nonzero monomial on $\Pi_i$, then the monomial is a square if and only if its restriction is a square.
\end{proof}

From now on we will assume that $\Delta = \Pi$ as in Section~\ref{sec-Pi}. Assume that $\Pi$ has vertices $v_0, v_1, \ldots, v_n$. The matrix $A=(\a{i, j})$ has size $n\times (n+1)$, with columns indexed by $0,1,\ldots,n$ and rows by $1,\ldots,n$. We use the notation $X_j$, $j=0,1,\ldots,n$, for the determinant of $A$ with its $j$-th column removed. If $J= (j_1, \ldots, j_s)$ is a vector with entries in $\{0,1,\ldots,n\}$, we let
\[ X_J = X_{j_1} X_{j_2} \cdots X_{j_s}.\]

Theorem~\ref{thm-PPconj} for $\Pi$ can be further reduced to the following:

\begin{theorem} \label{thm-simplest}
Let $I$ and $J$ be vectors with entries in $\{0,1,\dots,n\}$. Assume that $|I|=n$ and $|J|$ is odd.  Then
  \[ \partial_I  X_J  = 
  \begin{cases}
    (X_L)^2 & \text{if $X_I X_J = X_L^2 X_{(0,1,\dots,n)}$}, \\
    0         &  \text{otherwise.}
  \end{cases}\]
\end{theorem}

\begin{lemma} \label{lem-conj-Pi}
Theorem~\ref{thm-simplest} implies Theorem~\ref{thm-PPconj} for $\Pi$.
\end{lemma}
\begin{proof}
Let $I$ and $J$ be as in the statement of Theorem~\ref{thm-PPconj}, and let $J'$ be such that $X_{J'} = X_J X_{(0,1,\dots,n)}$. Note that $|J'|=2n+1$ is odd. We claim that Theorem~\ref{thm-PPconj} for $I,J$ is equivalent to Theorem~\ref{thm-simplest} for $I, J'$. 

Using that 
\[ W_\Pi (x_{J}) = \frac{X_{J}}{c_\Pi}  = \frac{X_{J}}{ X_{(0,1,\dots,n)}}, \]
the statement of Theorem~\ref{thm-PPconj} for $\Pi, I,J$ is
  \[ \partial_I  \frac{X_J}{c_\Pi}  = 
  \begin{cases}
    \frac{X_I X_J}{c_\Pi^2} & \text{if $\sqrt{x_I x_J}$ exists}, \\
    0         &  \text{otherwise.}
  \end{cases}\]
The statement of Theorem~\ref{thm-simplest} for $I, J'$ is 
  \[ \partial_I  ( X_J c_\Pi )  = 
  \begin{cases}
    X_I X_J & \text{if $\sqrt{X_I X_J}$ exists}, \\
    0         &  \text{otherwise.}
  \end{cases}\]
These two equations differ by a factor of $c_\Pi^2$.
\end{proof}
